\chapter{Foundations of Probability}
\section{Sample Space and $\sigma$-algebras}

\begin{definitionenv}
    The set of all possible outcomes of a random experiment is called the
    \textbf{sample space}, denoted by $\Omega$.
    Each element $\omega\in\Omega$ is called an \textbf{outcome}
    (or \textbf{sample point}).
\end{definitionenv}

\begin{exampleenv}
    \begin{itemize}
        \item Flip a coin once:
        $$
        \Omega=\{H,T\}.
        $$

        \item Flip a coin twice:
        $$
        \Omega=\{HH,HT,TH,TT\}.
        $$
    \end{itemize}
\end{exampleenv}

\begin{definitionenv}
    The collection $\mathcal{F}$ of subsets of $\Omega$ is a
    $\sigma$-algebra on $\Omega$ if
    \begin{enumerate}[label=(\alph*)]
        \item $\varnothing\in\mathcal{F}$.

        \item If
        $$
        A_1,A_2,\ldots,A_n,\ldots\in\mathcal{F},
        $$
        then
        $$
        \bigcup_{n}A_n\in\mathcal{F}.
        $$

        \item If $A\in\mathcal{F}$, then
        $$
        A^c\in\mathcal{F}.
        $$
    \end{enumerate}
    An \textbf{event} is a subset of the sample space $\Omega$.

    An event $A$ is called \textbf{measurable} if $A\in\mathcal{F}.$

    $(\Omega,\mathcal{F})$ is called \textbf{measure space}.
\end{definitionenv}

\begin{exampleenv}
    Let
    $$
    \Omega=\{HH,HT,TH,TT\}.
    $$
    Consider the following collections of subsets of $\Omega$:
    \begin{enumerate}[label=(\arabic*)]
        \item
        $$
        \mathcal{F}=\{\varnothing,\Omega\}.
        $$

        \item
        $$
        \mathcal{F}=2^{\Omega}.
        $$

        \item
        $$
        \mathcal{F}
        =\bigl\{\varnothing,\Omega,\{HH\},\{HT,TH,TT\}\bigr\}.
        $$

        \item
        $$
        \mathcal{F}
        =\bigl\{\varnothing,\Omega,\{HH,HT\},\{TH,TT\}\bigr\}.
        $$
    \end{enumerate}
    The $\sigma$-algebra specifies what information is available.
\end{exampleenv}

\begin{propositionenv}
    If $\mathcal{F}$ is a $\sigma$-algebra and
    $A_1,A_2,\ldots\in\mathcal{F}$, then
    $$
    \bigcap_{n=1}^{\infty} A_n\in\mathcal{F}.
    $$
\end{propositionenv}

\begin{proofenv}
    By De Morgan's law,
    $$
    \bigcap_{n=1}^{\infty} A_n
    =\left(\bigcup_{n=1}^{\infty} A_n^c\right)^c.
    $$
    Since $A_n\in\mathcal{F}$, we have $A_n^c\in\mathcal{F}$ for every $n$.
    Hence
    $$
    \bigcup_{n=1}^{\infty} A_n^c\in\mathcal{F},
    $$
    and therefore
    $$
    \left(\bigcup_{n=1}^{\infty} A_n^c\right)^c\in\mathcal{F}.
    $$
\end{proofenv}

\begin{exerciseenv}
    Let $\mathcal{F}$ be a $\sigma$-algebra on a finite sample space $\Omega$.
    How many elements can $\mathcal{F}$ contain?
\end{exerciseenv}

\begin{solutionenv}
    Let $|\Omega|=n$. Since $\Omega$ is finite, every $\sigma$-algebra
    $\mathcal{F}\subseteq 2^\Omega$ is finite.

    Let
    \[
    B_1,\ldots,B_k
    \]
    be the minimal nonempty elements of $\mathcal{F}$.
    These sets are pairwise disjoint and form a partition of $\Omega$.

    Moreover, every $A\in\mathcal{F}$ is a union of some of the
    $B_i$'s. Conversely, every union of the $B_i$'s belongs to
    $\mathcal{F}$.

    Hence
    \[
    \mathcal{F}
    =
    \left\{
        \bigcup_{i\in I} B_i:
        I\subseteq\{1,\ldots,k\}
    \right\}.
    \]

    There are $2^k$ subsets $I\subseteq\{1,\ldots,k\}$, so
    \[
    |\mathcal{F}|=2^k,
    \qquad 1\le k\le n.
    \]

    Therefore, the number of elements of a finite $\sigma$-algebra
    must be a power of $2$.
\end{solutionenv}

\begin{definitionenv}
    Given a measurable space $(\Omega,\mathcal{F})$, a \textbf{probability measure}
    on $(\Omega,\mathcal{F})$ is a function
    \[
        \begin{array}{rcl}
            \mathbb{P}:\mathcal{F} &\longrightarrow& [0,1]
        \end{array}
    \]
    satisfying:
    \begin{enumerate}[label=(\alph*)]
        \item
        \[
            \mathbb{P}(\Omega)=1.
        \]

        \item If $A_1,A_2,\ldots\in\mathcal{F}$ are pairwise disjoint, then
        \[
            \mathbb{P}\left(\bigcup_{n=1}^{\infty}A_n\right)
            =\sum_{n=1}^{\infty}\mathbb{P}(A_n).
        \]
    \end{enumerate}

    $(\Omega,\mathcal{F},\mathbb{P})$ is  called \textbf{probability space}.
\end{definitionenv}

\begin{propositionenv}
    Let $\mathbb{P}$ be a probability measure on $(\Omega,\mathcal{F})$. Then:
    \begin{enumerate}[label=(\arabic*)]
        \item $\mathbb{P}(\varnothing)=0$.
        \item For every $A\in\mathcal{F}$,
        \[
            \mathbb{P}(A^c)=1-\mathbb{P}(A).
        \]
        \item For every $A,B\in\mathcal{F}$,
        \[
            \mathbb{P}(A\cup B)=\mathbb{P}(A)+\mathbb{P}(B)-\mathbb{P}(A\cap B).
        \]
    \end{enumerate}
\end{propositionenv}

\begin{proofenv}
    \begin{enumerate}[label=(\arabic*)]
        \item Since $\varnothing$ and $\Omega$ are disjoint,
        \[
            \mathbb{P}(\varnothing\cup\Omega)=\mathbb{P}(\varnothing)+\mathbb{P}(\Omega)=\mathbb{P}(\Omega).
        \]
        , hence
        \[
            \mathbb{P}(\varnothing)=0.
        \]

        \item Since $A\cap A^c=\varnothing$ and $A\cup A^c=\Omega$,
        \[
            1=\mathbb{P}(\Omega)=\mathbb{P}(A\cup A^c)=\mathbb{P}(A)+\mathbb{P}(A^c),
        \]
        so
        \[
            \mathbb{P}(A^c)=1-\mathbb{P}(A).
        \]

        \item Write
        \[
            A\cup B=A\cup(B\setminus A),
        \]
        where $A$ and $B\setminus A$ are disjoint. Hence
        \[
            \mathbb{P}(A\cup B)=\mathbb{P}(A)+\mathbb{P}(B\setminus A).
        \]
        Also,
        \[
            B=(B\setminus A)\cup(A\cap B)
        \]
        is a disjoint union, so
        \[
            \mathbb{P}(B)=\mathbb{P}(B\setminus A)+\mathbb{P}(A\cap B).
        \]
        Therefore
        \[
            \mathbb{P}(A\cup B)=\mathbb{P}(A)+\mathbb{P}(B)-\mathbb{P}(A\cap B).
        \]
    \end{enumerate}
\end{proofenv}

\begin{propositionenv}
    Let $\mathbb{P}$ be a probability measure on $(\Omega,\mathcal{F})$.

    \begin{enumerate}[label=(\roman*)]
        \item If
        \[
            A_1 \subseteq A_2 \subseteq \cdots \subseteq A_n \subseteq \cdots,
        \]
        then the sequence $\{\mathbb{P}(A_n)\}_{n\ge 1}$ is non-decreasing and
        \[
            \lim_{n\to\infty}\mathbb{P}(A_n)
            =
            \mathbb{P}\left(\bigcup_{k=1}^{\infty} A_k\right).
        \]

        \item If
        \[
            B_1 \supseteq B_2 \supseteq \cdots \supseteq B_n \supseteq \cdots,
        \]
        then the sequence $\{\mathbb{P}(B_n)\}_{n\ge 1}$ is non-increasing and
        \[
            \lim_{n\to\infty}\mathbb{P}(B_n)
            =
            \mathbb{P}\left(\bigcap_{k=1}^{\infty} B_k\right).
        \]
    \end{enumerate}
\end{propositionenv}

\begin{proofenv}
    \begin{enumerate}[label=(\roman*)]
        \item Let
        \[
            A=\bigcup_{k=1}^{\infty}A_k
        \]
        and set $A_0=\varnothing$. Since $A_1\subseteq A_2\subseteq\cdots$,
        \[
            A
            =
            A_1
            \cup (A_2\setminus A_1)
            \cup (A_3\setminus A_2)
            \cup \cdots,
        \]
        where the sets on the right-hand side are pairwise disjoint. Hence
        \[
            \mathbb{P}(A)
            =
            \sum_{k=1}^{\infty}
            \mathbb{P}(A_k\setminus A_{k-1}).
        \]
        Therefore
        \[
        \begin{aligned}
            \mathbb{P}(A)
            &=
            \lim_{n\to\infty}
            \sum_{k=1}^{n}
            \mathbb{P}(A_k\setminus A_{k-1}) \\
            &=
            \lim_{n\to\infty}\mathbb{P}(A_n).
        \end{aligned}
        \]
        Thus
        \[
            \mathbb{P}(A_n)
            \uparrow
            \mathbb{P}\left(\bigcup_{k=1}^{\infty}A_k\right).
        \]

                \item Since
        \[
            B_1^c\subseteq B_2^c\subseteq\cdots,
        \]
        part (i) gives
        \[
            \mathbb{P}(B_n^c)
            \uparrow
            \mathbb{P}\left(\bigcup_{k=1}^{\infty}B_k^c\right).
        \]
        Using De Morgan's law and
        $\mathbb{P}(B_n)=1-\mathbb{P}(B_n^c)$,
        \[
        \begin{aligned}
            \lim_{n\to\infty}\mathbb{P}(B_n)
            &=
            1-\mathbb{P}\left(\bigcup_{k=1}^{\infty}B_k^c\right) \\
            &=
            \mathbb{P}\left(\bigcap_{k=1}^{\infty}B_k\right).
        \end{aligned}
        \]
        Hence
        \[
            \mathbb{P}(B_n)
            \downarrow
            \mathbb{P}\left(\bigcap_{k=1}^{\infty}B_k\right).
        \]
    \end{enumerate}
\end{proofenv}

\begin{definitionenv}
    For a real sequence $(a_n)_{n\ge1}$, define
    \[
        \limsup_{n\to\infty} a_n
        :=
        \inf_{n\ge1}\sup_{k\ge n} a_k,
        \qquad
        \liminf_{n\to\infty} a_n
        :=
        \sup_{n\ge1}\inf_{k\ge n} a_k.
    \]
    Equivalently, they are the limits of the decreasing tail suprema
    $\sup_{k\ge n}a_k$ and the increasing tail infima
    $\inf_{k\ge n}a_k$, respectively.
\end{definitionenv}

\begin{definitionenv}
    For a sequence of sets $(A_n)_{n\ge1}$, define
    \[
        \limsup_{n\to\infty} A_n
        :=
        \bigcap_{n=1}^{\infty}\bigcup_{k\ge n} A_k,
        \qquad
        \liminf_{n\to\infty} A_n
        :=
        \bigcup_{n=1}^{\infty}\bigcap_{k\ge n} A_k.
    \]
    Equivalently,
    \[
        x\in\limsup_{n\to\infty}A_n
        \iff
        x\in A_n \text{ for infinitely many }n,
    \]
    while
    \[
        x\in\liminf_{n\to\infty}A_n
        \iff
        x\in A_n \text{ for all but finitely many }n.
    \]
    In particular,
    \[
        \liminf_{n\to\infty}A_n
        \subseteq
        \limsup_{n\to\infty}A_n.
    \]
\end{definitionenv}


\begin{definitionenv}
    Let $A,B\in\mathcal{F}$ with $\mathbb{P}(B)>0$.
    The \textbf{conditional probability} of $A$ given $B$ is defined by
    \[
        \mathbb{P}(A\mid B)
        :=
        \frac{\mathbb{P}(A\cap B)}{\mathbb{P}(B)}.
    \]
\end{definitionenv}

\begin{propositionenv}
    Let $B_1,B_2,\ldots$ be pairwise disjoint events such that
    \[
        \bigcup_k B_k=\Omega
    \]
    and
    \[
        \mathbb{P}(B_k)>0
        \qquad\text{for every }k.
    \]
    Then, for every event $A$,
    \[
        \mathbb{P}(A)
        =
        \sum_k
        \mathbb{P}(A\mid B_k)\mathbb{P}(B_k).
    \]
\end{propositionenv}

\begin{proofenv}
    Since the sets $B_k$ are pairwise disjoint and
    $\bigcup_k B_k=\Omega$,
    \[
        A
        =
        A\cap\Omega
        =
        A\cap\left(\bigcup_k B_k\right)
        =
        \bigcup_k (A\cap B_k),
    \]
    where the sets $A\cap B_k$ are pairwise disjoint. Hence
    \[
        \mathbb{P}(A)
        =
        \sum_k \mathbb{P}(A\cap B_k).
    \]
    By the definition of conditional probability,
    \[
        \mathbb{P}(A\cap B_k)
        =
        \mathbb{P}(A\mid B_k)\mathbb{P}(B_k).
    \]
    Therefore
    \[
        \mathbb{P}(A)
        =
        \sum_k
        \mathbb{P}(A\mid B_k)\mathbb{P}(B_k).
    \]
\end{proofenv}

\begin{definitionenv}
    Two events $A$ and $B$ are said to be \textbf{independent} if
    \[
        \mathbb{P}(A\cap B)
        =
        \mathbb{P}(A)\mathbb{P}(B).
    \]

    More generally, a sequence of events
    \[
        A_1,A_2,\ldots
    \]
    is said to be \textbf{independent} if, for every finite collection of
    distinct indices $k_1,\ldots,k_m$,
    \[
        \mathbb{P}\left(
            \bigcap_{j=1}^{m} A_{k_j}
        \right)
        =
        \prod_{j=1}^{m}
        \mathbb{P}(A_{k_j}).
    \]
\end{definitionenv}

\begin{remark}
    For events $(A_k)_{k\ge1}$,
    \[
        \limsup_{k\to\infty}A_k
    \]
    is the event that infinitely many of the $A_k$'s occur.
\end{remark}

\begin{propositionenv}
    Let $(A_k)_{k\ge 1}$ be a sequence of events.

    \begin{enumerate}[label=(\arabic*)]
        \item If
        \[
            \sum_{k=1}^{\infty}\mathbb{P}(A_k)<\infty,
        \]
        then
        \[
            \mathbb{P}\left(
                \limsup_{k\to\infty}A_k
            \right)=0.
        \]

        \item If
        \[
            \sum_{k=1}^{\infty}\mathbb{P}(A_k)=\infty
        \]
        and the events $A_1,A_2,\ldots$ are independent, then
        \[
            \mathbb{P}\left(
                \limsup_{k\to\infty}A_k
            \right)=1.
        \]
    \end{enumerate}
\end{propositionenv}

\begin{proofenv}
    \begin{enumerate}[label=(\arabic*)]
        \item For every $n$,
        \[
            \limsup_{k\to\infty}A_k
            \subseteq
            \bigcup_{k\ge n}A_k.
        \]
        Hence
        \[
            0
            \le
            \mathbb{P}\left(\limsup_{k\to\infty}A_k\right)
            \le
            \mathbb{P}\left(\bigcup_{k\ge n}A_k\right)
            \le
            \sum_{k=n}^{\infty}\mathbb{P}(A_k).
        \]
        Letting $n\to\infty$, the tail sum tends to $0$, so
        \[
            \mathbb{P}\left(\limsup_{k\to\infty}A_k\right)=0.
        \]

        \item For $m\ge n$, by independence,
        \[
            \mathbb{P}\left(\bigcap_{k=n}^{m}A_k^c\right)
            =
            \prod_{k=n}^{m}\bigl(1-\mathbb{P}(A_k)\bigr)
            \le
            \exp\left(
                -\sum_{k=n}^{m}\mathbb{P}(A_k)
            \right).
        \]
        Letting $m\to\infty$ gives
        \[
            \mathbb{P}\left(\bigcap_{k=n}^{\infty}A_k^c\right)=0.
        \]
        Therefore
        \[
        \begin{aligned}
            \mathbb{P}\left(
                \left(\limsup_{k\to\infty}A_k\right)^c
            \right)
            &=
            \mathbb{P}\left(
                \bigcup_{n=1}^{\infty}
                \bigcap_{k=n}^{\infty}A_k^c
            \right) \\
            &\le
            \sum_{n=1}^{\infty}
            \mathbb{P}\left(
                \bigcap_{k=n}^{\infty}A_k^c
            \right)
            =0.
        \end{aligned}
        \]
        Hence
        \[
            \mathbb{P}\left(\limsup_{k\to\infty}A_k\right)=1.
        \]
    \end{enumerate}
\end{proofenv}

\begin{exampleenv}[Gambler's ruin]
    Consider a random walk on the states
    \[
        \{0,1,\ldots,N\}.
    \]
    Starting from a state $k\in\{1,\ldots,N-1\}$, the walk moves
    \[
        k\to k+1 \quad \text{with probability } p,
    \]
    and
    \[
        k\to k-1 \quad \text{with probability } q,
    \]
    where
    \[
        p+q=1.
    \]

    Let
    \[
        \rho_k
        :=
        \mathbb{P}
        \bigl(
            \text{hit $N$ before $0$}
            \,\big|\,
            \text{start at $k$}
        \bigr).
    \]
    Then
    \[
        \rho_0=0,
        \qquad
        \rho_N=1.
    \]

    Conditioning on the first step gives, for $1\le k\le N-1$,
    \[
        \rho_k
        =
        p\rho_{k+1}+q\rho_{k-1}.
    \]
    Hence
    \[
        p(\rho_{k+1}-\rho_k)
        =
        q(\rho_k-\rho_{k-1}),
    \]
    so
    \[
        \rho_{k+1}-\rho_k
        =
        \frac{q}{p}
        (\rho_k-\rho_{k-1}).
    \]
    Therefore,
    \[
        \rho_k-\rho_{k-1}
        =
        \left(\frac{q}{p}\right)^{k-1}
        (\rho_1-\rho_0).
    \]

    Since
    \[
        1
        =
        \rho_N-\rho_0
        =
        \sum_{k=1}^{N}
        (\rho_k-\rho_{k-1}),
    \]
    we obtain
    \[
        1
        =
        \left(
            1+\frac{q}{p}
            +\left(\frac{q}{p}\right)^2
            +\cdots
            +\left(\frac{q}{p}\right)^{N-1}
        \right)
        (\rho_1-\rho_0).
    \]

    \begin{enumerate}[label=(\arabic*)]
        \item If
        \[
            p=q=\frac12,
        \]
        then
        \[
            \rho_1-\rho_0=\frac1N,
        \]
        and therefore
        \[
            \rho_k
            =
            \sum_{j=1}^{k}
            (\rho_j-\rho_{j-1})
            =
            \frac{k}{N}.
        \]

        \item If $p\ne q$, then
        \[
            \rho_1-\rho_0
            =
            \frac{1-\frac{q}{p}}
                 {1-\left(\frac{q}{p}\right)^N}.
        \]
        Hence
        \[
            \rho_k-\rho_{k-1}
            =
            \left(\frac{q}{p}\right)^{k-1}
            \frac{1-\frac{q}{p}}
                 {1-\left(\frac{q}{p}\right)^N},
        \]
        and summing from $1$ to $k$ yields
        \[
            \rho_k
            =
            \frac{
                1-\left(\frac{q}{p}\right)^k
            }{
                1-\left(\frac{q}{p}\right)^N
            }.
        \]
    \end{enumerate}

    Thus
    \[
        \boxed{
        \rho_k
        =
        \begin{cases}
            \dfrac{k}{N},
            & p=q=\dfrac12, \\[1.2ex]
            \dfrac{
                1-\left(\frac{q}{p}\right)^k
            }{
                1-\left(\frac{q}{p}\right)^N
            },
            & p\ne q.
        \end{cases}}
    \]
\end{exampleenv}

\begin{definitionenv}
A \textbf{random variable} is a function
\[
    X:\Omega\to\mathbb R
\]
such that, for every open set $A\subseteq\mathbb R$,
\[
    X^{-1}(A)
    =
    \{\omega\in\Omega:X(\omega)\in A\}
    \in\mathcal F.
\]
Equivalently, $X$ is $\mathcal F$-measurable.
\end{definitionenv}

\begin{exampleenv}

Let
\[
    \Omega=\{HH,HT,TH,TT\}.
\]

Define $X_1,X_2:\Omega\to\mathbb R$ by
\[
    X_1(HH)=X_1(HT)=1,
    \qquad
    X_1(TH)=X_1(TT)=-1,
\]
and
\[
    X_2(HH)=X_2(TH)=1,
    \qquad
    X_2(HT)=X_2(TT)=-1.
\]

Let
\[
    X=X_1+X_2.
\]

Then the smallest $\sigma$-algebras that make $X_1$, $X_2$, and $X$
random variables are, respectively,
\[
    \mathcal F_1
    =
    \sigma\bigl(\{HH,HT\},\{TH,TT\}\bigr),
\]
\[
    \mathcal F_2
    =
    \sigma\bigl(\{HH,TH\},\{HT,TT\}\bigr),
\]
and
\[
    \mathcal F_3
    =
    \sigma\bigl(\{HH\},\{HT,TH\},\{TT\}\bigr).
\]

\end{exampleenv}

\begin{exampleenv}
Let
\[
    \Omega=\{1,2,3,4,5,6\},
    \qquad
    \mathcal F
    =
    \sigma\bigl(\{1,3,5\},\{2,4,6\}\bigr).
\]

The function
\[
    X_1(k)=-k
\]
is not $\mathcal F$-measurable, while
\[
    X_2(k)
    =
    \begin{cases}
        1, & k \text{ is odd},\\
        2, & k \text{ is even},
    \end{cases}
\]
is $\mathcal F$-measurable.
\end{exampleenv}

\begin{definitionenv}
\begin{enumerate}
    \item Let $p\in[0,1]$. A random variable $X$ has the
\textbf{Bernoulli distribution} with parameter $p$, written
\[
    X\sim\operatorname{Ber}(p),
\]
if
\[
    \mathbb{P}(X=1)=p,
    \qquad
    \mathbb{P}(X=0)=1-p.
\]

\item Let $n\in\mathbb{N}_+$ and $p\in[0,1]$. A random variable $X$ has the
\textbf{binomial distribution} with parameters $(n,p)$, written
\[
    X\sim\operatorname{Bin}(n,p),
\]
if
\[
    \mathbb{P}(X=k)
    =
    \binom{n}{k}p^k(1-p)^{n-k},
    \qquad
    k=0,\ldots,n.
\]
\end{enumerate}
\end{definitionenv}

\begin{remark}
If $p$ is small and $n$ is large with
\[
    np=\lambda>0,
\]
then, for fixed $k$,
\[
\begin{aligned}
    \mathbb{P}(X=k)
    &=
    \binom{n}{k}
    \left(\frac{\lambda}{n}\right)^k
    \left(1-\frac{\lambda}{n}\right)^{n-k} \\
    &\longrightarrow
    \frac{\lambda^k}{k!}e^{-\lambda}.
\end{aligned}
\]
\end{remark}

\begin{definitionenv}
\begin{enumerate}
    \item Let $\lambda>0$. A random variable $X$ has the
\textbf{Poisson distribution} with parameter $\lambda$, written
\[
    X\sim\operatorname{Pois}(\lambda),
\]
if
\[
    \mathbb{P}(X=k)
    =
    e^{-\lambda}\frac{\lambda^k}{k!},
    \qquad
    k=0,1,2,\ldots.
\]

\item Let $p\in(0,1]$. A random variable $X$ has the
\textbf{geometric distribution} with parameter $p$, written
\[
    X\sim\operatorname{Geo}(p),
\]
if
\[
    \mathbb{P}(X=k)
    =
    (1-p)^{k-1}p,
    \qquad
    k\ge1.
\]
\end{enumerate}
\end{definitionenv}

\begin{propositionenv}
If
\[
    X\sim\operatorname{Geo}(p),
\]
then
\[
    \mathbb{P}(X>n)=(1-p)^n,
    \qquad
    n\ge0.
\]
\end{propositionenv}

\begin{proofenv}
\[
\begin{aligned}
    \mathbb{P}(X>n)
    &=
    \sum_{k=n+1}^{\infty}
    p(1-p)^{k-1} \\
    &=
    p(1-p)^n
    \sum_{k=0}^{\infty}(1-p)^k \\
    &=
    (1-p)^n.
\end{aligned}
\]
\end{proofenv}

\begin{propositionenv}
A geometric random variable has the \textbf{memoryless property}:
\[
    \mathbb{P}(X>m+n\mid X>m)
    =
    \mathbb{P}(X>n).
\]
\end{propositionenv}

\begin{proofenv}
\[
\begin{aligned}
    \mathbb{P}(X>m+n\mid X>m)
    &=
    \frac{\mathbb{P}(X>m+n)}
         {\mathbb{P}(X>m)} \\
    &=
    \frac{(1-p)^{m+n}}{(1-p)^m} \\
    &=
    (1-p)^n \\
    &=
    \mathbb{P}(X>n).
\end{aligned}
\]
\end{proofenv}

\begin{exerciseenv}
Let $X$ be a discrete random variable satisfying
\[
    \mathbb{P}(X>m+n\mid X>m)
    =
    \mathbb{P}(X>n)
\]
for all $m,n\ge0$. Show that
\[
    X\sim\operatorname{Geo}(p)
\]
for some $p\in(0,1]$.
\end{exerciseenv}

\begin{remark}[Stirling's formula]
\[
    n!
    \sim
    \sqrt{2\pi n}
    \left(\frac ne\right)^n.
\]
That is,
\[
    \lim_{n\to\infty}
    \frac{n!}
    {\sqrt{2\pi n}\left(\frac ne\right)^n}
    =1.
\]

Taking logarithms in Stirling's formula gives
\[
\begin{aligned}
    \log(n!)
    &=
    \frac12\log(2\pi n)
    +n(\log n-1)
    +o(1).
\end{aligned}
\]
Hence,
\[
    \lim_{n\to\infty}
    \frac{\log(n!)}{n\log n}
    =1.
\]
\end{remark}

\begin{definitionenv}
Let $X:\Omega\to\mathbb R$ be a random variable. The
\textbf{distribution function} of $X$ is
\[
    F:\mathbb R\to[0,1],
\]
defined by
\[
    F(x)=\mathbb P(X\le x).
\]

If $F$ is differentiable, then
\[
    f(x):=F'(x)
\]
is called the \textbf{probability density function} of $X$.
\end{definitionenv}

\begin{propositionenv}
The distribution function $F$ is right continuous.
\end{propositionenv}

\begin{proofenv}
Let
\[
    x_n\downarrow x.
\]
Then
\[
    \{X\le x_1\}\supseteq
    \{X\le x_2\}\supseteq\cdots
\]
and hence
\[
\begin{aligned}
    \lim_{n\to\infty}F(x_n)
    &=
    \lim_{n\to\infty}\mathbb P(X\le x_n)\\
    &=
    \mathbb P\left(
        \bigcap_{n=1}^{\infty}\{X\le x_n\}
    \right)\\
    &=
    \mathbb P(X\le x)\\
    &=F(x).
\end{aligned}
\]
Therefore $F$ is right continuous.
\end{proofenv}

\begin{exampleenv}
If
\[
    X\sim\operatorname{Ber}(p),
\]
then
\[
    F(x)=\mathbb P(X\le x)
\]
is given by
\[
F(x)=
\begin{cases}
0,&x<0,\\
1-p,&0\le x<1,\\
1,&x\ge1.
\end{cases}
\]
\end{exampleenv}

\begin{remark}
For a continuous random variable $X$ with density function $f$,
\[
    \mathbb P(a<X<b)
    =
    F(b)-F(a)
    =
    \int_a^b f(x)\,dx.
\]
\end{remark}

\begin{exampleenv}
A random variable $X$ is uniformly distributed on $[a,b]$, written
\[
    X\sim\operatorname{Unif}[a,b],
\]
if its density function is
\[
    f(x)=
    \begin{cases}
        \dfrac1{b-a},&x\in[a,b],\\[6pt]
        0,&\text{otherwise}.
    \end{cases}
\]
\end{exampleenv}

\begin{definitionenv}
Let $\lambda>0$. A random variable $X$ has the
\textbf{exponential distribution} with parameter $\lambda$, written
\[
    X\sim\operatorname{Exp}(\lambda),
\]
if its density function is
\[
f(x)=
\begin{cases}
    \lambda e^{-\lambda x},&x\ge0,\\
    0,&x<0.
\end{cases}
\]
\end{definitionenv}


\begin{propositionenv}
If
\[
    X\sim\operatorname{Exp}(\lambda),
\]
then
\[
    \mathbb P(X>t)=e^{-\lambda t},
    \qquad t\ge0.
\]
\end{propositionenv}

\begin{proofenv}
\[
\begin{aligned}
    \mathbb P(X>t)
    &=
    \int_t^\infty
    \lambda e^{-\lambda x}\,dx\\
    &=
    \left[-e^{-\lambda x}\right]_t^\infty\\
    &=
    e^{-\lambda t}.
\end{aligned}
\]
\end{proofenv}


\begin{propositionenv}
The exponential distribution satisfies the memoryless property:
\[
    \mathbb P(X>s+t\mid X>s)
    =
    \mathbb P(X>t).
\]
\end{propositionenv}

\begin{proofenv}
\[
\begin{aligned}
    \mathbb P(X>s+t\mid X>s)
    &=
    \frac{\mathbb P(X>s+t)}
    {\mathbb P(X>s)}\\
    &=
    \frac{e^{-\lambda(s+t)}}
    {e^{-\lambda s}}\\
    &=
    e^{-\lambda t}\\
    &=
    \mathbb P(X>t).
\end{aligned}
\]
\end{proofenv}

\begin{exerciseenv}
Show that the exponential distribution is the only continuous distribution
satisfying the memoryless property.
\end{exerciseenv}

\begin{definitionenv}
A random variable $X$ has the normal distribution with parameters
$\mu\in\mathbb R$ and $\sigma>0$, written
\[
    X\sim N(\mu,\sigma^2),
\]
if its density function is
\[
    f(x)
    =
    \frac1{\sqrt{2\pi}\sigma}
    \exp\left(
        -\frac{(x-\mu)^2}{2\sigma^2}
    \right),
    \qquad x\in\mathbb R.
\]
\end{definitionenv}

\begin{remark}
The constant in the density is chosen so that
\[
    \int_{\mathbb R}f(x)\,dx=1.
\]
It remains to show that
\[
    \int_{\mathbb R}e^{-x^2/2}\,dx=\sqrt{2\pi}.
\]
\end{remark}

\begin{proofenv}
Let
\[
    I=\int_{\mathbb R}e^{-x^2/2}\,dx.
\]
Then
\[
\begin{aligned}
    I^2
    &=
    \iint_{\mathbb R^2}
    e^{-(x^2+y^2)/2}\,dxdy .
\end{aligned}
\]

Using polar coordinates,
\[
    x=r\cos\theta,\qquad y=r\sin\theta,
\]
we obtain
\[
\begin{aligned}
    I^2
    &=
    \int_0^{2\pi}
    \int_0^\infty
    e^{-r^2/2}r\,drd\theta\\
    &=
    \int_0^{2\pi}
    \int_0^\infty
    e^{-u}\,du\,d\theta
    \qquad
    \left(u=\frac{r^2}{2}\right)\\
    &=2\pi .
\end{aligned}
\]
Hence
\[
    I=\sqrt{2\pi}.
\]
\end{proofenv}

\begin{definitionenv}
Let $X$ be a random variable.

If $X$ is discrete and
\[
    \mathbb P(X=x_k)=p_k,
\]
with
\[
    \sum_k p_k|x_k|<+\infty,
\]
then the expectation of $X$ is defined by
\[
    \mathbb EX
    :=
    \sum_k p_kx_k.
\]
\end{definitionenv}


\begin{definitionenv}
If $X$ has density function $f$ and
\[
    \int_{\mathbb R}|x|f(x)\,dx<+\infty,
\]
then
\[
    \mathbb EX
    :=
    \int_{\mathbb R}x f(x)\,dx.
\]
\end{definitionenv}


\begin{remark}
More generally, on a probability space $(\Omega,\mathcal F,\mathbb P)$,
if
\[
    \int_\Omega |X|\,d\mathbb P<+\infty,
\]
then
\[
    \mathbb EX
    =
    \int_\Omega X\,d\mathbb P.
\]
\end{remark}

\begin{definitionenv}
The variance of a random variable $X$ is defined by
\[
    \operatorname{Var}(X)
    :=
    \mathbb E(X-\mathbb EX)^2.
\]
Equivalently,
\[
    \operatorname{Var}(X)
    =
    \mathbb EX^2-(\mathbb EX)^2.
\]
\end{definitionenv}


\begin{theoremenv}
\textbf{(Jensen's inequality)}
Let
\[
    J:\mathbb R\to\mathbb R
\]
be convex. Then
\[
    \mathbb E J(X)
    \ge
    J(\mathbb EX).
\]
\end{theoremenv}


\begin{exampleenv}
For $p\ge1$,
\[
    J(x)=|x|^p
\]
is convex. Hence
\[
    \mathbb E|X|^p
    \ge
    |\mathbb EX|^p.
\]
\end{exampleenv}


\begin{theoremenv}
\textbf{(Hölder's inequality)}
Let $p,q\in[1,\infty)$ satisfy
\[
    \frac1p+\frac1q=1.
\]
Then
\[
    |\mathbb E(XY)|
    \le
    (\mathbb E|X|^p)^{1/p}
    (\mathbb E|Y|^q)^{1/q}.
\]
\end{theoremenv}